% 原书物理358–361／印刷335–338；69条原提示，不增添解答.
\clearpage\phantomsection\section*{提示}\label{sec:hints}
\addcontentsline{toc}{section}{提示}
\begin{sourcehint}{1}
使用分段仿射的振荡函数. 对 $C^1$ 函数不可能做到，因为像测度具有有界密度的 $C^1$ 函数 $f$ 应满足 $f'\neq0$，并且是微分同胚. 但在微分同胚中，只有 $f(x)=x$ 和 $f(x)=1-x$ 保测度.
\end{sourcehint}
\begin{sourcehint}{2}
写出 Kantorovich 问题，并在目标测度中使用变量替换 $y\mapsto-y$.
\end{sourcehint}
\begin{sourcehint}{3}
寻找一个单调径向运输. 答案是
\[
\mathrm T(x)=4\sqrt{1-\sqrt{1-|x|^2}}\,\frac{x}{|x|}.
\]
\end{sourcehint}
\begin{sourcehint}{4}
对两个点使用 $c$-循环单调性（即单调性）.
\end{sourcehint}
\begin{sourcehint}{5}
对第一个问题，使用凸函数梯度的刻画. 对第二个问题，再次使用 $c$-循环单调性.
\end{sourcehint}
\begin{sourcehint}{6}
对第一个费用，使用 Jensen 不等式. 对第二个费用，将它写成 $(x_1+\cdots+x_N)^2$ 加上分别只依赖各个 $x_i$ 的函数.
\end{sourcehint}
\begin{sourcehint}{7}
可以（部分地）使用费用为 $h(y-x)$ 的运输问题，其中 $z<0$ 时 $h(z)=+\infty$，$z\geq0$ 时 $h(z)=0$.
\end{sourcehint}
\begin{sourcehint}{8}
将 $\gamma$ 分解为对角线上方与下方两部分，并使用一维凸费用的结果. 然后证明，两者重叠并非最优.
\end{sourcehint}
\begin{sourcehint}{9}
以函数 $\phi(x)=2-4|x-\frac12|$ 和 $\psi(y)=2-4|y-\frac12|$ 作为 Kantorovich 势的猜测，验证以下定义的 $\mathrm T$ 最优：当 $x<\frac12$ 时，$\mathrm T(x)=x+\frac12$；当 $x\geq\frac12$ 时，$\mathrm T(x)=x-\frac12$.
\end{sourcehint}
\begin{sourcehint}{10}
若存在两对 $(x_0,y_0),(x_1,y_1)\in\spt\gamma$，使 $x_0<x_1$ 且 $y_0>y_1$，在 $[x_0,x_1]\times[y_1,y_0]$ 上积分 $\frac{\partial^2c}{\partial x\partial y}$，得到矛盾.
\end{sourcehint}
\begin{sourcehint}{11}
将 $A^{-1}B$ 写成一个对称正定矩阵与一个旋转的乘积.
\end{sourcehint}
\begin{sourcehint}{12}
$u(x)=-x$ 是一个势.
\end{sourcehint}
\begin{sourcehint}{13}
思路与\autoref{ex:3}相同. 用势 $u(x)=-|x|$ 证明所找到的映射最优.
\end{sourcehint}
\begin{sourcehint}{14}
费用 $c(x,y)=|x-y|^\alpha$ 是一个距离.
\end{sourcehint}
\begin{sourcehint}{15}
使用与\autoref{sec:3.3.2}相同的分析.
\end{sourcehint}
\begin{sourcehint}{16}
使用\autoref{sec:3.3.1}的分解策略.
\end{sourcehint}
\begin{sourcehint}{17}
使用关于 $L^\infty$ 最优运输的结果.
\end{sourcehint}
\begin{sourcehint}{18}
费用 $c_\varepsilon$ 一致收敛于线性费用.
\end{sourcehint}
\begin{sourcehint}{19}
考虑费用关于 $\varepsilon$ 的 Taylor 展开，以猜测低阶项.
\end{sourcehint}
\begin{sourcehint}{20}
从将有理数集增厚的经典例子出发.
\end{sourcehint}
\begin{sourcehint}{21}
证明，所有满足 $|u(x)|\leq P(x)$ 的函数 $u:X\to\R$ 都属于 SNCF 距离下的 $\Lip_1$，然后应用对偶公式，得到 $\mathscr T_c$ 的下界. 对于相反方向的界，只需让公共质量保持静止.
\end{sourcehint}
\begin{sourcehint}{22}
只需用 $C_c^1$ 函数 $\psi(x)$ 检验.
\end{sourcehint}
\begin{sourcehint}{23}
两种方法：对维数归纳；使用多项式卷积（采用类似 $x^n(1-x)^n$ 的函数）.
\end{sourcehint}
\begin{sourcehint}{24}
由唯一性，$\sigma$ 是径向的. Kantorovich 势为 $u(x)=-|x|$. 使用 Monge--Kantorovich 方程.
\end{sourcehint}
\begin{sourcehint}{25}
使用与\autoref{obj:4.24}相同的策略. 另见\cite{b273}.
\end{sourcehint}
\begin{sourcehint}{26}
这等价于存在 $\mathbf w\in L^p$，满足 $\nabla\cdot\mathbf w=\mu-\nu$.
\end{sourcehint}
\begin{sourcehint}{27}
求解极小化问题
\[
\min\{i_{\widetilde Q}(\Omega):(e_0)_\#\widetilde Q=\mu,\ (e_1)_\#\widetilde Q=\nu,\ i_{\widetilde Q}\leq i_Q\}.
\]
\end{sourcehint}
\begin{sourcehint}{28}
作如下设想：从 $A$ 到 $A'$ 的质量分成三部分，$\alpha$ 经过 $B$，$\alpha$ 经过 $B'$，$\beta$ 直接到 $A'$；从 $B$ 到 $B'$ 也同样处理. 写出均衡条件对应的一个 $2\times2$ 方程组.
\end{sourcehint}
\begin{sourcehint}{29}
在 $\Omega$ 内的一个立方体中取规则网格. 计算将每个小立方体中央一半的质量移出该小立方体所需的费用，并求和.
\end{sourcehint}
\begin{sourcehint}{30}
使用角度定律\eqref{eq:4.31}.
\end{sourcehint}
\begin{sourcehint}{31}
假设 $\phi\geq0$，使用关于下半连续函数的\autoref{box:1.5}中定义的 $k$-Lipschitz 函数 $\phi_k$.
\end{sourcehint}
\begin{sourcehint}{32}
证明，最优 $\nu$ 是 $f$ 通过到 $\Sigma$ 的投影所得的像. 使用 Beckmann 表述.
\end{sourcehint}
\begin{sourcehint}{33}
用\autoref{obj:5.5}修改 $\gamma$ 的边缘.
\end{sourcehint}
\begin{sourcehint}{34}
对 $W_2^2$ 求导.
\end{sourcehint}
\begin{sourcehint}{35}
证明 $\Gamma$-收敛.
\end{sourcehint}
\begin{sourcehint}{36}
取达到 $\max d(x,\spt(\nu))$ 的点 $x_0\in\spt(\mu)$：它必须被运输到某处.
\end{sourcehint}
\begin{sourcehint}{37}
构造一个反例：映射 $\mathrm T$ 是梯度，但对 $h(z)=|z|^4$，$\nabla h(\mathrm T(x)-x)$ 不是梯度.
\end{sourcehint}
\begin{sourcehint}{38}
改用\autoref{obj:5.33}的证明.
\end{sourcehint}
\begin{sourcehint}{39}
移动一个点时，写出最优性条件.
\end{sourcehint}
\begin{sourcehint}{40}
考虑一个极点矩阵中属于 $]0,1[$ 的元素. 证明，若存在这样的元素，就存在一个“循环”，并得出矛盾.
\end{sourcehint}
\begin{sourcehint}{41}
使用与上面相同的策略.
\end{sourcehint}
\begin{sourcehint}{42}
由 $\log\gamma_{ij}^\varepsilon=\log p_i^\varepsilon+\log q_j^\varepsilon+\log\eta_{ij}$，推出 $\varepsilon\log p_i^\varepsilon+\varepsilon\log q_j^\varepsilon\leq c_{ij}$；若在极限下 $\gamma_{ij}>0$，则取等.
\end{sourcehint}
\begin{sourcehint}{43}
写出 $\cof(D^2u)D^2u=\det(D^2u)\mathrm I$ 并求导.
\end{sourcehint}
\begin{sourcehint}{44}
计算二阶导数并分部积分，得到 $\int u_{ij}v^iv^jf$.
\end{sourcehint}
\begin{sourcehint}{45}
第一部分：加上再减去适当的、关于 $W_1$ 连续的泛函，使被积函数非负，并将其限制到有界区域上. 第二部分：证明二阶矩的界蕴含 $W_1$ 收敛.
\end{sourcehint}
\begin{sourcehint}{46}
使用以下事实：若起点与终点测度的 $L^\infty$ 范数都被1控制，则沿测地线的 $L^\infty$ 范数也被1控制.
\end{sourcehint}
\begin{sourcehint}{47}
用直接法证明存在性，用严格凸性证明唯一性，并在最优性条件中引入势.
\end{sourcehint}
\begin{sourcehint}{48}
较难：证明一种内球条件，即若 $\varrho(x_0)>g(x_0)$，则在 $B(\mathrm T(x_0),|x_0-\mathrm T(x_0)|)$ 上 $\varrho=f$，并使用 Lebesgue 点. 见\cite{b151}.
\end{sourcehint}
\begin{sourcehint}{49}
使用 $\int\frac12|\nabla\varrho|^2+4\sqrt\varrho$.
\end{sourcehint}
\begin{sourcehint}{50}
考虑 $\varrho_0$ 与 $\varrho$ 之间的最优 $\gamma$：若 $\varrho$ 的原子多于 $N$ 个，则将每个 $\gamma\restr(\{x_i^0\}\times\R^d)$ 替换为其重心，可以得到更好的结果.
\end{sourcehint}
\begin{sourcehint}{51}
证明重心保持不变，并用方程计算 $E'$.
\end{sourcehint}
\begin{sourcehint}{52}
用势写出最优性条件（用直接法证明存在性，用严格凸性证明\mbox{唯一性）.}
\end{sourcehint}
\begin{sourcehint}{53}
用 $\|\varrho\|_{L^m}^2$ 估计 $\nabla u_\varrho$ 部分，并利用 $m>2$ 说明这些范数沿极小化序列必须有界.
\end{sourcehint}
\begin{sourcehint}{54}
为约束引入 Lagrange 乘子，将问题写成极小极大问题，并推出存在鞍点的条件.
\end{sourcehint}
\begin{sourcehint}{55}
写出并求解 Euler--Lagrange 方程.
\end{sourcehint}
\begin{sourcehint}{56}
用 Young 不等式和\autoref{obj:5.33}，证明 $t\mapsto\|\varrho_t-1\|_{H^{-1}}^2$ 属于 BV.
\end{sourcehint}
\begin{sourcehint}{57}
验证，可以将 $\nu$ 的每个原子移到 $\nu$ 的最近的（某个）原子上.
\end{sourcehint}
\begin{sourcehint}{58}
使用凹费用及具有公共质量的测度.
\end{sourcehint}
\begin{sourcehint}{59}
尝试实现费用0.
\end{sourcehint}
\begin{sourcehint}{60}
采用\autoref{sec:1.4}反例中的 $\mu$ 与 $\nu$，并将 $c(x,y)$ 选为 $|x-y|$ 的一个不连续函数.
\end{sourcehint}
\begin{sourcehint}{61}
改造映满正方形的 Peano 构造，将映射定义为某个压缩算子的不动点.
\end{sourcehint}
\begin{sourcehint}{62}
使用最优映射为凸函数梯度的刻画（二次费用）. 在 $\R$ 上两个紧支集测度之间找出一个非最优运输，并将它们平移到相距很远的位置（线性费用）. 找到一个不是梯度的两个梯度的复合.
\end{sourcehint}
\begin{sourcehint}{63}
使用 $F_\mu^{[-1]}$ 的性质.
\end{sourcehint}
\begin{sourcehint}{64}
将最优运输写成 $x-\varphi'(x)$，验证 $\|1-g\|_{\dot H^{-1}}=\|\varphi'\|_{L^2}$，并分部积分.
\end{sourcehint}
\begin{sourcehint}{65}
对于 $\mathscr T_c\geq0$，用方案 $\gamma=\mu\otimes\nu$ 得到一个界. 对于 $\gamma$ 的支集，取 $(x,y)$ 和 $(x',y')$ 属于 $\spt(\gamma)$，写出 $c$-循环单调性，再对 $(x',y')$ 积分. 这给出了一个上界，可用于证明上半连续性.
\end{sourcehint}
\begin{sourcehint}{66}
利用熵的位移凸性，比较沿一条测地线的熵在 $t=0$ 与 $t=1$ 处的导数.
\end{sourcehint}
\begin{sourcehint}{67}
对于第一个问题，对 $W_2^2$ 求导并使用\autoref{ex:66}. 对于第二个问题，使用\autoref{obj:5.2}中的不等式及热方程的半群性质.
\end{sourcehint}
\begin{sourcehint}{68}
对 $W_2^2$ 求导后，使用 Gronwall 引理. $H^1$ 是合理的，因为 $H^1$ 范数在时间上受控.
\end{sourcehint}
\begin{sourcehint}{69}
对演化方程的解与静态解之间的 $W_2^2$ 求导.
\end{sourcehint}
